\section{Conclusion}
\label{sec:conclusion}

We asked whether there are tasks where classical ML struggles while quantum
methods---simulated at $20$ qubits or more---excel, and we answered it with a
single reproducible study that demonstrates both sides with honest scope. On an
engineered projected-kernel task at $20$ qubits, a quantum-kernel SVM reaches
$0.865$ test accuracy while the best of six tuned classical models reaches $0.629$
($p=3.9\times10^{-8}$, $d=5.4$), and the advantage is stable at $\approx
0.22$--$0.23$ across $8$, $10$, and $20$ qubits. On the discrete-log task,
off-the-shelf classical ML stays at chance across $12$/$17$/$21$-bit problems
while the quantum interval-state kernel separates the task near-perfectly
($0.997 \to 1.000$).

\para{Key takeaway.} A $20{+}$-qubit simulated quantum kernel demonstrably
generalizes on tasks where the best classical models cannot---but the win is the
\emph{right} quantum inductive bias, not quantumness alone: a generic fidelity
kernel concentrates to chance exactly like a classical overfitter. The honest
scope is engineered or quantum-structured tasks under noiseless simulation; a
naturally labeled simulable advantage remains open.

\para{Future work.} We see five concrete next steps: (i) repeat Experiment~A under
finite-shot kernel estimation to test whether the advantage survives $O(2^{-d})$
eigenvalue concentration; (ii) sweep $N$ and bandwidth to map the generalization
curve versus sample size; (iii) add a classical-surrogate baseline (neural tangent
kernel) as the strictest comparison; (iv) replace random inputs with PCA-reduced
Fashion-MNIST to reproduce the advantage on real-image features; and (v) explore
quantum-\emph{data} tasks, such as learning properties of many-body states, where
advantage is least contested.
